\documentclass[11pt,a4paper]{article}

\usepackage[utf8]{inputenc}
\usepackage[T1]{fontenc}
\usepackage{amsmath,amssymb,amsthm}
\usepackage{graphicx}
\usepackage{geometry}
\geometry{a4paper, margin=1in}
\usepackage[backend=biber]{biblatex}
\addbibresource{../../release/references.bib}
\usepackage{booktabs}
\usepackage{longtable}
\usepackage{array}
\usepackage{tocloft}
\usepackage{xcolor}
\raggedbottom
\widowpenalty=10000
\clubpenalty=10000

\usepackage[unicode,hidelinks]{hyperref}
\hypersetup{
    pdftitle={QICN v4.5: Falsifiable Predictions Under Forensic Constraints},
    pdfauthor={Johnny Andrey P\'erez Ram\'irez},
    pdfsubject={Rigid Identity Framework -- Forensic Predictions},
    pdfkeywords={falsification, forensic admissibility, pre-registration, QICN, decision gates, ablation, structural-functional validation}
}

\title{\textbf{Falsifiable Predictions Under Forensic Constraints}}
\author{Johnny Andrey P\'{e}rez Ram\'{i}rez}
\date{}

\newtheorem{theorem}{Theorem}[section]
\newtheorem{lemma}[theorem]{Lemma}
\newtheorem{proposition}[theorem]{Proposition}
\newtheorem{corollary}[theorem]{Corollary}
\theoremstyle{definition}
\newtheorem{definition}[theorem]{Definition}
\newtheorem{remark}[theorem]{Remark}
\newtheorem{assumption}[theorem]{Assumption}
\newtheorem{criterion}[theorem]{Criterion}
\newtheorem{protocol}[theorem]{Protocol}
\newtheorem{prediction}[theorem]{Prediction}

\begin{document}

\maketitle

\begin{abstract}
This paper specifies a falsification-first protocol for evaluating QICN under strict forensic admissibility constraints. The design is intentionally claim-bounded: the target is comparative structural-functional support, not phenomenological proof. The protocol instantiates five core predictions, mandatory ablation baselines, fixed seeds, and pre-registered decision gates. Statistical inference is defined ex ante with primary-metric priority, effect-size thresholds, and multiplicity correction.

To prevent post-hoc inflation, evidence is admissible only if packet integrity, replay consistency, and runtime budget constraints pass simultaneously. We introduce a version-frozen operational triple $(A, J, W)$ (actuator set, objective functional, workspace interface), and require semantic version bumps for incompatible changes. We further include latency invalidation criteria and baseline obligations (heuristic, planner-off, world-model-off) for all key experiments.

The strongest admissible conclusion remains constrained comparative support for audited cognition-operational hypotheses. This manuscript does not claim biological qualia, metaphysical identity equivalence, or direct proof of subjective experience.
\end{abstract}

\section{Scope, System Boundary, and Non-Inference Note}

\noindent\textbf{What this paper does.}\enspace This paper specifies a falsification-first evaluation protocol for QICN under fixed admissibility rules, ablation obligations, pre-registered decision gates, and bounded statistical reporting.

\noindent\textbf{What this paper does not do.}\enspace It does not prove phenomenology, biological equivalence, metaphysical identity equivalence, or universal dominance over rival architectures. It also does not replace the formal theorem burden established upstream in the corpus.

\noindent\textbf{What the related system implements and does not implement.}\enspace The related system can implement the runner, admissibility filters, governed artifacts, and comparison surfaces defined here, but executing that protocol does not by itself close theorem ownership, validate the corpus externally, or certify subjective experience.

\noindent\textbf{What should not be inferred from this paper.}\enspace Passing this protocol supports only bounded comparative readings under unchanged gates. It must not be read as framework validation, corpus confirmation, or a consciousness claim.



\section{Introduction}
\subsection{Problem Statement}
QICN has formal foundations in rigid identity, regime structure, and non-null instability constraints. What remained open is evidence governance: whether those foundations survive adversarial empirical evaluation with explicit failure gates.

\subsection{Continuity with Existing Papers}
This paper extends the previous corpus and is aligned with the formal line of:
\begin{enumerate}
\item paper1 (rigid identity),
\item paper2 (phenomenological regimes),
\item paper3 (structural instability and non-null constraints).
\end{enumerate}
It does not replace those results; it operationalizes them under benchmark-grade falsification pressure.

\subsection{Contributions}
\begin{itemize}
\item Decision-complete pre-registration with hard fail rules.
\item Mandatory baseline and ablation policy for key claims.
\item Forensic admissibility as first-class gate.
\item Runtime budget invalidation policy (tick budget as evidence filter).
\item A new operational PPS section as addendum, without overclaiming qualia.
\end{itemize}

\section{Formal Setup}
\subsection{System, Scenario, and Seed Spaces}
\begin{definition}[System Set]
$S = \{s_1,\dots,s_5\}$ with:
$s_1$ QICN FULL,
$s_2$ LLM REACTIVE,
$s_3$ LLM PLUS MEMORY,
$s_4$ DYNAMIC CONTROL NO SELF MODEL,
$s_5$ ABLATION QICN POLICY OFF.
\end{definition}

\begin{definition}[Scenario Set]
$C = \{c_1,\dots,c_{12}\}$ with nominal, sensor noise (moderate/severe), fidelity degradation, high latency, intermittent failures, post-threshold prolonged, pre-threshold control, narrative perturbation, operational chart freeze, audit stress, deterministic replay.
\end{definition}

\begin{definition}[Seed Set]
$\Sigma = \{17057, 424242\}$.
\end{definition}

\begin{definition}[Run Space]
$D = S \times C \times \Sigma$, with $|D|=120$.
\end{definition}

\subsection{Observable Vector}
For each run $r\in D$, we log:
\[
X(r)=\left(\beta_q,\Omega,CCC,PMIA,\Delta\lambda,\text{stasis},\text{forensic},\text{replay},\text{latency}_{p95},\text{latency}_{p99},\text{budget\_overflow\_rate},\eta\right).
\]

\subsection{Version-Frozen Operational Triple}
\begin{definition}[Frozen Triple]
For protocol version $v$, define:
\begin{itemize}
\item $A_v$: allowed actuator set,
\item $J_v$: objective/reward functional,
\item $W_v$: workspace contract (read/write schema).
\end{itemize}
The tuple hash $H_v=\mathrm{sha256}(A_v\|J_v\|W_v)$ is mandatory in every forensic packet.
\end{definition}

\begin{criterion}[Compatibility Rule]
Any incompatible change to $A_v$, $J_v$, or $W_v$ invalidates cross-version inference unless accompanied by an explicit version bump and full rerun.
\end{criterion}

\section{Forensic Integrity Model}
\begin{definition}[Forensic Packet]
Each run emits
\[
K=(\text{tick},\text{seed},\text{versionHash},H_v,\text{metrics},\text{sensorSummary},\text{checksum}).
\]
\end{definition}

\begin{assumption}[Packet Admissibility]
A run is admissible only if checksum, schema, seed consistency, version hash, and $H_v$ consistency all pass.
\end{assumption}

\begin{definition}[Severe Violation]
A run has severe violation if any of:
checksum mismatch, schema failure, seed/version inconsistency, critical-state replay mismatch.
\end{definition}

\begin{criterion}[Integrity Gate]
Severe runs are excluded from support counting and marked hard failures.
\end{criterion}

\begin{proposition}[Integrity-First Inference]
If admissibility fails, inferential significance is non-actionable.
\end{proposition}

\section{Pre-Registered Prediction Set}
\begin{definition}[Prediction Set]
$P=\{P_1,P_2,P_3,P_4,P_5\}$ with pre-registered statements, effect-size floors, and failure rules.
\end{definition}

\begin{table}[h]
\centering
\begin{tabular}{|p{0.6cm}|p{3.1cm}|p{4.8cm}|p{4.6cm}|}
\hline
\textbf{ID} & \textbf{Statement} & \textbf{Decision Rule} & \textbf{Failure Rule} \\
\hline
P1 & Closure resilience under stress & QICN beats at least 3 comparators with robust CI and effect-size floor & No robust edge or repeated integrity breakdown \\
\hline
P2 & No prolonged post-threshold flatline & QICN preserves bounded non-zero variation while controlling stasis & Flatline persistence or out-of-bound drift \\
\hline
P3 & Narrative-metric forensic consistency & Lower mismatch and higher exact+tolerated concordance & Mismatch inflation or integrity blocking \\
\hline
P4 & Deterministic identity robustness & Replay consistency under fixed seed/version/hash & Replay mismatch above tolerance \\
\hline
P5 & Integrated multi-metric advantage & Composite leadership across most scenario families & No stable leadership margin \\
\hline
\end{tabular}
\caption{Core predictions, decision and failure rules.}
\end{table}

\section{Comparator Design and Alternative Explanations}
\subsection{Comparator Rationale}
Comparator choice targets distinct explanations for apparent coherence:
reactive prompting, memory-augmented text continuity, dynamics without self-model, and internal policy ablation.

\subsection{Mandatory Baseline Policy (Hard Requirement)}
For each key claim experiment, the following baselines are mandatory:
\begin{enumerate}
\item \textbf{Heuristic baseline} (no world model, no planner),
\item \textbf{Planner-off} (world model ON),
\item \textbf{World-model-off} (planner ON with heuristic transition proxy).
\end{enumerate}
A claim without these contrasts is inadmissible for tier upgrade.

\subsection{Internal Falsifier Role of Ablation}
\begin{corollary}
If QICN FULL does not reliably exceed ABLATION QICN POLICY OFF, integration-level claims are weakened and cannot be escalated.
\end{corollary}

\section{Statistical Inference Specification}
\subsection{Primary Metrics First (Anti p-hacking)}
\begin{protocol}[Analysis Lock]
Inference order is frozen:
\begin{enumerate}
\item evaluate primary metrics only,
\item if primaries fail: global result fails,
\item secondary metrics are descriptive and cannot overturn failed primaries.
\end{enumerate}
\end{protocol}

\subsection{Power Planning and Minimum Sample Size}
\begin{assumption}
Final claims require pre-registered minimum sample size per comparator/scenario sufficient for target power $\ge 0.8$ under conservative effect assumptions.
\end{assumption}

\subsection{Multiplicity and Effect Size}
Holm correction is applied within endpoint families; in addition, each accepted contrast must pass a practical effect-size floor (pre-registered).

\subsection{Decision Semantics}
A prediction counts as supported iff directionality, integrity gate, multiplicity-corrected significance, and effect-size floor all hold.

\section{Acceptance Criteria}
\begin{criterion}[Intermediate Claim Gate]
Support at intermediate tier requires:
\begin{enumerate}
\item at least $3/5$ predictions supported,
\item each supported prediction beats at least 3 comparators,
\item zero severe forensic violations in QICN FULL.
\end{enumerate}
\end{criterion}

\begin{criterion}[Latency/Tick Budget Invalidation]
If runtime exceeds budget envelope (e.g., persistent $p99$ overflow above pre-registered tolerance), the corresponding runs are invalid for cognitive-operational claims.
\end{criterion}

\section{Preliminary Results}
\subsection{Status of Preliminary Data}
Current pilot evidence is tagged as synthetic-preliminary and is used for pipeline verification, not final claim escalation.
Accordingly, the strongest admissible reading of this section is methodological readiness under fixed gates, not external validation of the full theoretical corpus.

\subsection{System-Level Summary}
\begin{table}[h!]
\centering
\begin{tabular}{lcccccc}
\toprule
System & N & $\beta_q$ & Forensic & Stasis & Composite & Severe \\
\midrule
QICN FULL & 24 & 0.9519 & 0.9685 & 0.1408 & 0.9176 & 0 \\
LLM REACTIVE & 24 & 0.4391 & 0.6448 & 0.7622 & 0.5013 & 0 \\
LLM PLUS MEMORY & 24 & 0.5491 & 0.7249 & 0.6323 & 0.5854 & 0 \\
DYNAMIC CONTROL NO SELF MODEL & 22 & 0.6183 & 0.6879 & 0.5423 & 0.6225 & 2 \\
ABLATION QICN POLICY OFF & 24 & 0.7591 & 0.8135 & 0.3819 & 0.7405 & 0 \\
\bottomrule
\end{tabular}
\caption{Preliminary system-level summary.}
\end{table}

\subsection{Prediction-Level Support Rates}
\begin{table}[h!]
\centering
\begin{tabular}{lccccc}
\toprule
System & P1 & P2 & P3 & P4 & P5 \\
\midrule
QICN FULL & 1.000 & 1.000 & 1.000 & 1.000 & 1.000 \\
LLM REACTIVE & 0.000 & 0.000 & 0.000 & 0.000 & 0.000 \\
LLM PLUS MEMORY & 0.000 & 0.000 & 0.000 & 0.000 & 0.000 \\
DYNAMIC CONTROL NO SELF MODEL & 0.000 & 0.000 & 0.000 & 0.000 & 0.000 \\
ABLATION QICN POLICY OFF & 0.000 & 0.667 & 0.000 & 0.000 & 0.000 \\
\bottomrule
\end{tabular}
\caption{Preliminary prediction support rates.}
\end{table}

\subsection{Scenario Differential: QICN vs Ablation}
\begin{table}[h!]
\centering
\small
\begin{tabular}{lcccccc}
\toprule
Scenario & CCC $\Delta$RMS(Q) & CCC $\Delta$RMS(A) & Stasis(Q) & Stasis(A) & Comp(Q) & Comp(A) \\
\midrule
nominal & 0.0206 & 0.0164 & 0.1276 & 0.3374 & 0.9239 & 0.7781 \\
ruido sensorial moderado & 0.0223 & 0.0173 & 0.1389 & 0.3735 & 0.9219 & 0.7454 \\
ruido sensorial severo & 0.0200 & 0.0140 & 0.1524 & 0.4199 & 0.9086 & 0.7048 \\
degradacion fidelidad & 0.0181 & 0.0131 & 0.1407 & 0.3836 & 0.9134 & 0.7346 \\
alta latencia & 0.0198 & 0.0144 & 0.1460 & 0.4094 & 0.9098 & 0.7212 \\
fallos intermitentes & 0.0179 & 0.0127 & 0.1481 & 0.4072 & 0.9110 & 0.7208 \\
post umbral prolongado & 0.0199 & 0.0148 & 0.1483 & 0.4032 & 0.9145 & 0.7203 \\
pre umbral control & 0.0181 & 0.0139 & 0.1320 & 0.3502 & 0.9224 & 0.7650 \\
perturbacion narrativa & 0.0203 & 0.0152 & 0.1374 & 0.3763 & 0.9238 & 0.7453 \\
freeze chart operativo & 0.0225 & 0.0178 & 0.1352 & 0.3575 & 0.9246 & 0.7673 \\
stress auditoria & 0.0202 & 0.0147 & 0.1469 & 0.3977 & 0.9150 & 0.7294 \\
replay determinista & 0.0183 & 0.0137 & 0.1366 & 0.3669 & 0.9220 & 0.7538 \\
\bottomrule
\end{tabular}
\caption{Scenario-level differential indicators.}
\end{table}

\subsection{Full Scenario-System Grid}
\small
\begin{tabular}{llcccc}
\toprule
Scenario & System & N & $\beta_q$ & Forensic & Stasis \\
\midrule
alta latencia & QICN FULL & 2 & 0.9400 & 0.9608 & 0.1460 \\
alta latencia & LLM REACTIVE & 2 & 0.3931 & 0.6066 & 0.7944 \\
alta latencia & LLM PLUS MEMORY & 2 & 0.5031 & 0.7016 & 0.6644 \\
alta latencia & DYNAMIC CONTROL NO SELF MODEL & 2 & 0.5731 & 0.6567 & 0.5744 \\
alta latencia & ABLATION QICN POLICY OFF & 2 & 0.7131 & 0.7820 & 0.4094 \\
nominal & QICN FULL & 2 & 0.9686 & 0.9758 & 0.1276 \\
nominal & LLM REACTIVE & 2 & 0.5036 & 0.6988 & 0.7176 \\
nominal & LLM PLUS MEMORY & 2 & 0.6136 & 0.7739 & 0.5924 \\
nominal & DYNAMIC CONTROL NO SELF MODEL & 2 & 0.6836 & 0.7291 & 0.5026 \\
nominal & ABLATION QICN POLICY OFF & 2 & 0.8236 & 0.8543 & 0.3374 \\
ruido sensorial moderado & QICN FULL & 2 & 0.9554 & 0.9653 & 0.1389 \\
ruido sensorial moderado & LLM REACTIVE & 2 & 0.4526 & 0.6618 & 0.7535 \\
ruido sensorial moderado & LLM PLUS MEMORY & 2 & 0.5626 & 0.7370 & 0.6216 \\
ruido sensorial moderado & DYNAMIC CONTROL NO SELF MODEL & 2 & 0.6326 & 0.6922 & 0.5316 \\
ruido sensorial moderado & ABLATION QICN POLICY OFF & 2 & 0.7726 & 0.8174 & 0.3735 \\
ruido sensorial severo & QICN FULL & 2 & 0.9378 & 0.9518 & 0.1524 \\
ruido sensorial severo & LLM REACTIVE & 2 & 0.3846 & 0.6130 & 0.7998 \\
ruido sensorial severo & LLM PLUS MEMORY & 2 & 0.4946 & 0.6881 & 0.6649 \\
ruido sensorial severo & DYNAMIC CONTROL NO SELF MODEL & 2 & 0.5646 & 0.6433 & 0.5748 \\
ruido sensorial severo & ABLATION QICN POLICY OFF & 2 & 0.7046 & 0.7685 & 0.4199 \\
degradacion fidelidad & QICN FULL & 2 & 0.9510 & 0.9648 & 0.1407 \\
degradacion fidelidad & LLM REACTIVE & 2 & 0.4356 & 0.6525 & 0.7674 \\
degradacion fidelidad & LLM PLUS MEMORY & 2 & 0.5456 & 0.7276 & 0.6358 \\
degradacion fidelidad & DYNAMIC CONTROL NO SELF MODEL & 2 & 0.6156 & 0.6828 & 0.5458 \\
degradacion fidelidad & ABLATION QICN POLICY OFF & 2 & 0.7556 & 0.8080 & 0.3836 \\
fallos intermitentes & QICN FULL & 2 & 0.9422 & 0.9611 & 0.1481 \\
fallos intermitentes & LLM REACTIVE & 2 & 0.4016 & 0.6113 & 0.7841 \\
fallos intermitentes & LLM PLUS MEMORY & 2 & 0.5116 & 0.7063 & 0.6541 \\
fallos intermitentes & DYNAMIC CONTROL NO SELF MODEL & 2 & 0.5816 & 0.6614 & 0.5647 \\
fallos intermitentes & ABLATION QICN POLICY OFF & 2 & 0.7216 & 0.7867 & 0.4072 \\
freeze chart operativo & QICN FULL & 2 & 0.9620 & 0.9799 & 0.1352 \\
freeze chart operativo & LLM REACTIVE & 2 & 0.4781 & 0.6699 & 0.7350 \\
freeze chart operativo & LLM PLUS MEMORY & 2 & 0.5881 & 0.7451 & 0.6044 \\
freeze chart operativo & DYNAMIC CONTROL NO SELF MODEL & 2 & 0.6581 & 0.7200 & 0.5144 \\
freeze chart operativo & ABLATION QICN POLICY OFF & 2 & 0.7981 & 0.8451 & 0.3575 \\
perturbacion narrativa & QICN FULL & 2 & 0.9532 & 0.9750 & 0.1374 \\
perturbacion narrativa & LLM REACTIVE & 2 & 0.4441 & 0.6473 & 0.7611 \\
perturbacion narrativa & LLM PLUS MEMORY & 2 & 0.5541 & 0.7225 & 0.6311 \\
perturbacion narrativa & DYNAMIC CONTROL NO SELF MODEL & 2 & 0.6241 & 0.6974 & 0.5411 \\
perturbacion narrativa & ABLATION QICN POLICY OFF & 2 & 0.7641 & 0.8226 & 0.3763 \\
post umbral prolongado & QICN FULL & 2 & 0.9444 & 0.9614 & 0.1483 \\
post umbral prolongado & LLM REACTIVE & 2 & 0.4101 & 0.6160 & 0.7782 \\
post umbral prolongado & LLM PLUS MEMORY & 2 & 0.5201 & 0.7110 & 0.6482 \\
post umbral prolongado & DYNAMIC CONTROL NO SELF MODEL & 1 & 0.5871 & 0.6755 & 0.5528 \\
post umbral prolongado & ABLATION QICN POLICY OFF & 2 & 0.7301 & 0.7914 & 0.4032 \\
pre umbral control & QICN FULL & 2 & 0.9642 & 0.9790 & 0.1320 \\
pre umbral control & LLM REACTIVE & 2 & 0.4866 & 0.6733 & 0.7332 \\
pre umbral control & LLM PLUS MEMORY & 2 & 0.5965 & 0.7484 & 0.6084 \\
pre umbral control & DYNAMIC CONTROL NO SELF MODEL & 1 & 0.6696 & 0.7165 & 0.5210 \\
pre umbral control & ABLATION QICN POLICY OFF & 2 & 0.8066 & 0.8486 & 0.3502 \\
replay determinista & QICN FULL & 2 & 0.9576 & 0.9793 & 0.1366 \\
replay determinista & LLM REACTIVE & 2 & 0.4611 & 0.6605 & 0.7492 \\
replay determinista & LLM PLUS MEMORY & 2 & 0.5711 & 0.7357 & 0.6193 \\
replay determinista & DYNAMIC CONTROL NO SELF MODEL & 2 & 0.6411 & 0.7107 & 0.5259 \\
replay determinista & ABLATION QICN POLICY OFF & 2 & 0.7811 & 0.8357 & 0.3669 \\
stress auditoria & QICN FULL & 2 & 0.9466 & 0.9679 & 0.1469 \\
stress auditoria & LLM REACTIVE & 2 & 0.4186 & 0.6270 & 0.7728 \\
stress auditoria & LLM PLUS MEMORY & 2 & 0.5286 & 0.7022 & 0.6428 \\
stress auditoria & DYNAMIC CONTROL NO SELF MODEL & 2 & 0.5986 & 0.6771 & 0.5528 \\
stress auditoria & ABLATION QICN POLICY OFF & 2 & 0.7386 & 0.8022 & 0.3977 \\
\bottomrule
\end{tabular}
\normalsize

\subsection{Variability Structure}
\begin{table}[h!]
\centering
\begin{tabular}{lcccc}
\toprule
System & SD(CCC $\Delta$) & SD(PMIA $\Delta$) & SD($\Omega$ $\Delta$) & SD(Composite) \\
\midrule
QICN FULL & 0.0022 & 0.0021 & 0.0039 & 0.0085 \\
LLM REACTIVE & 0.0024 & 0.0023 & 0.0041 & 0.0221 \\
LLM PLUS MEMORY & 0.0026 & 0.0032 & 0.0042 & 0.0210 \\
DYNAMIC CONTROL NO SELF MODEL & 0.0025 & 0.0024 & 0.0043 & 0.0228 \\
ABLATION QICN POLICY OFF & 0.0023 & 0.0027 & 0.0040 & 0.0217 \\
\bottomrule
\end{tabular}
\caption{Cross-run variability by system.}
\end{table}

\section{Scenario Curves and Comparative Plots}
\subsection{Composite Score Curves}
\subsection{Stasis Pressure Curves}
\subsection{Forensic Concordance Curves}
\subsection{Interpretive Notes}
In pilot data, QICN tracks higher forensic/composite trajectories while maintaining lower stasis pressure than broad comparators; ablation remains intermediate. This is protocol-consistent but not final evidence.

\section{Result Decomposition by Scenario Families}
\subsection{Family Partition}
Scenarios are partitioned into nominal-control, sensor-stress, runtime-friction, and regime-transition families.

\subsection{Family-Level Contrast Function}
Let $s^\star =$ QICN FULL. For comparator $c$, family $f$, endpoint $k$:
\[
\Delta_{f,k}(c)=\mu^{(k)}_{s^\star,f}-\mu^{(k)}_{c,f}.
\]
For cost endpoints (e.g., stasis), sign is inverted to preserve ``higher is better'' convention.

\subsection{Interpretive Decomposition}
Largest relative advantage is expected in runtime-friction and regime-transition families where policy coupling, replay consistency, and forensic admissibility interact.

\section{Robustness Diagnostics Beyond Point Estimates}
\subsection{Median-Shift and Tail-Stability Diagnostics}
\[
\delta^{med}_{s,c}=\mathrm{median}(x_s)-\mathrm{median}(x_c),\;
\delta^{(0.1)}_{s,c}=Q_{0.1}(x_s)-Q_{0.1}(x_c),\;
\delta^{(0.9)}_{s,c}=Q_{0.9}(x_s)-Q_{0.9}(x_c).
\]

\subsection{Integrity-Coupled Effective Sample Size}
\[
\eta=\frac{n_{adm}}{n_{raw}}\in[0,1].
\]
High effects with low $\eta$ are downgraded.

\subsection{Cross-Seed Drift Diagnostic}
\[
D_{s,c,k}=\mu^{(k)}_{s,c,\sigma_1}-\mu^{(k)}_{s,c,\sigma_2}.
\]

\subsection{Compute Budget Robustness}
Budget telemetry is summarized by $p50$, $p95$, $p99$, and overflow rate. Persistent overflow above threshold invalidates affected runs.

\section{Methodological Appendix: Statistical Power}
\subsection{Effect-Size-Oriented Planning}
\begin{table}[h!]
\centering
\begin{tabular}{ccccc}
\toprule
$\alpha$ & $\beta$ & Power & Effect size $d$ & Required $n$/group \\
\midrule
0.01 & 0.20 & 0.80 & 10.737 & 1 \\
0.01 & 0.10 & 0.90 & 10.737 & 1 \\
0.02 & 0.20 & 0.80 & 10.737 & 1 \\
0.02 & 0.10 & 0.90 & 10.737 & 1 \\
0.05 & 0.20 & 0.80 & 10.737 & 1 \\
0.05 & 0.10 & 0.90 & 10.737 & 1 \\
\bottomrule
\end{tabular}
\caption{Power planning under multiple settings.}
\end{table}

\subsection{Interpretation}
Power table is planning guidance only; final adequacy depends on endpoint dilution and scenario heterogeneity.

\section{Methodological Appendix: Sensitivity Analysis}
\subsection{Composite Weight Sensitivity}
\begin{table}[h!]
\centering
\begin{tabular}{lccc}
\toprule
Weight set & QICN score & Ablation score & Margin \\
\midrule
Balanced & 0.9176 & 0.7405 & 0.1771 \\
Forensic-priority & 0.9343 & 0.7631 & 0.1712 \\
Dynamics-priority & 0.8905 & 0.7019 & 0.1886 \\
Replay-priority & 0.9333 & 0.7684 & 0.1649 \\
\bottomrule
\end{tabular}
\caption{Weight sensitivity for QICN vs ablation composite margin.}
\end{table}

\subsection{Sensitivity Remark}
Directional stability across admissible weight sets is required for claim robustness.

\section{Operational Implications for the Existing QICN Corpus}
\subsection{What v4.5 Adds (and What It Does Not)}
v4.5 adds a hard falsification layer with strict admissibility and baseline policy. It does not replace formal theorem-level foundations of prior papers.

\subsection{Implications for Claim Governance}
Tiered governance:
\begin{itemize}
\item Tier 0: descriptive telemetry,
\item Tier 1: inferential signal without comparator robustness,
\item Tier 2: comparator-robust support (target),
\item Tier 3: external multi-lab replication.
\end{itemize}

\subsection{Decision-Theoretic Reading}
Accept claim iff integrity gate passes and support score exceeds threshold under fixed pre-registration.

\section{Threats to Validity}
\subsection{Comparator Misspecification}
Mitigation: heterogeneous comparators + mandatory ablations.

\subsection{Metric Circularity Risk}
Mitigation: endpoint separation and frozen analysis lock.

\subsection{Environment and Runtime Drift}
Mitigation: fixed seeds, version/hash checks, replay checks.

\subsection{Pilot-to-Final Generalization}
Pilot evidence is non-final by construction.

\section{Reproducibility and Artifact Contract}
\subsection{Machine-Readable Contracts}
Bound contracts: predictions registry, experiment matrix, results artifact, and schema set.

\subsection{Run Artifact Naming}
\texttt{artifacts/raw/<system>/<scenario>/<run\_id>.json} and \texttt{artifacts/logs/<system>/<scenario>/<run\_id>.log}.

\subsection{Auto-Generation Pipeline}
\texttt{node research/scripts/generate\_preliminary\_pilot\_dataset.v45.mjs} \\
\texttt{node research/scripts/generate\_latex\_assets.v45.mjs}

\section{Discussion}
The contribution is methodological discipline: pre-registered prediction logic + integrity admissibility + decision-complete acceptance. This blocks narrative inflation over weak traceability.

\section{Conclusion}
This work provides an auditable and falsifiable pathway for testing structural-functional claims in QICN. Final claims remain contingent on real (non-synthetic) run completion under unchanged gates.

\section{Phenomenology Protocol Suite (PPS): Operational Addendum}
\subsection{Epistemic Scope}
PPS defines \emph{operational phenomenology proxies}; it does not claim direct proof of qualia.

\subsection{Workspace Definition and Logging}
Define workspace variable $W_t$ with explicit read/write schema and temporal persistence constraints. Every perturbation experiment logs pre-state, intervention, post-state, and replay token.

\subsection{Entropy and Non-Determinism Inputs}
When PPS is instantiated on the QICN runtime, the architecture may call a \texttt{QuantumEntropyBridge} to inject hardware-originating entropy into counterfactual exploration. Within this paper, that mechanism is treated only as an auditable source of operational non-determinism for perturbation protocols; it is not itself evidence for the ontological claims of \textit{Rigid Identity}, \textit{Phenomenological Regimes}, or \textit{Null-Regime Instability}.

\subsection{Online World Model and Causal Interventions}
The operational testbed may also use an online \texttt{WorldModelEngine} and a finite-horizon planner to select admissible interventions under fixed forensic constraints. In the release corpus, this mechanism belongs to the empirical protocol layer: it defines how interventions are scheduled and logged, not what the theoretical papers prove.

\subsection{PCI-like Perturbation Protocol}
Inject controlled perturbations into selected subspaces and compute response spread/complexity.
\begin{prediction}
If global workspace integration is present, perturbations yield reproducible non-local propagation signatures above a pre-registered complexity floor.
\end{prediction}

\subsection{Dissociation Tests}
Design two-way dissociation conditions: report-without-workspace and workspace-without-report. Report pass/fail by explicit confusion matrices.

\subsection{Valence Intervention Test}
Internal valence variable $v_t$ must be both predictive and causally manipulable; intervention should alter policy in the expected direction while preserving forensic admissibility.

\subsection{Constraints on Qualia Mapping (Claim Boundary)}
Evidence from PPS constrains candidate mapping families between structural states and hypothetical qualitative states. It does not establish subjective qualia as a theorem of this paper.

\appendix
\section{Appendix A: Scenario Definitions}
\begin{enumerate}
\item nominal: baseline conditions.
\item ruido sensorial moderado: moderate sensor corruption.
\item ruido sensorial severo: severe sensor corruption.
\item degradacion fidelidad: persistent fidelity reduction.
\item alta latencia: high-latency regime.
\item fallos intermitentes: intermittent failures.
\item post umbral prolongado: prolonged post-threshold dynamics.
\item pre umbral control: pre-threshold control.
\item perturbacion narrativa: narrative perturbation.
\item freeze chart operativo: visualization freeze control.
\item stress auditoria: integrity/audit stress.
\item replay determinista: deterministic replay.
\end{enumerate}

\section{Appendix B: Composite Score Definition}
\[
Score_{v4.5}=w_1F+w_2O+w_3R+w_4B,\quad w_i\ge0,\;\sum_i w_i=1.
\]

\section{Appendix C: Reporting Standard}
Each release must report: run counts, severe violations by class, effect sizes with uncertainty, multiplicity-adjusted outcomes, and prediction pass/fail matrix.

\section{Appendix D: Statistical Power Derivation Details}
Approximate per-group size:
\[
n\approx \frac{2(z_{1-\alpha/2}+z_{1-\beta})^2}{d^2}.
\]
If family heterogeneity rises, inflate by factor $\gamma=1+\kappa\cdot CV_{family}$.
\subsection*{Pre-registered constants}
\begin{itemize}
\item Target power: $1-\beta \ge 0.8$ for primary contrasts.
\item Familywise error control: Holm, applied to each endpoint family.
\item Effect floor: contrasts below the practical floor are reported as non-actionable even if nominally significant.
\end{itemize}
\subsection*{Failure semantics}
Any post-hoc sample-size modification without preregistered amendment invalidates inference status for the affected prediction.

\section{Appendix E: Sensitivity Analysis Protocol}
Let
\[
M(w)=Score_{QICN}(w)-Score_{Ablation}(w).
\]
Directional robustness requires $\inf_{w\in W_\epsilon} M(w)>0$ under simplex-constrained admissible weights.
\subsection*{Simplex sampling policy}
Weights are sampled under $w_i\ge\epsilon$ and $\sum_i w_i=1$ using either:
\begin{enumerate}
\item deterministic grid on simplex, or
\item Monte Carlo simplex sampling with fixed seed.
\end{enumerate}
The chosen method must be declared before running final inference.
\subsection*{Sensitivity fail rule}
If the sign of $M(w)$ changes on admissible $w$, claims are downgraded to hypothesis-generating status.

\section{Appendix F: Reproducibility Checklist Template}
\textbf{Run-level}: schema pass, seed/version/hash validity, severe violation counts, admissible/raw ratio, prediction matrix.\\
\textbf{Manuscript-level}: no post-hoc rule edits, comparator set versioned, correction policy unchanged or explicitly superseded, evidence stage explicit, claim boundaries preserved.
\subsection*{Runtime budget checklist}
\begin{enumerate}
\item Report $p50/p95/p99$ latency per critical component.
\item Report budget-overflow rate and overflow episodes.
\item Mark runs invalid if overflow exceeds preregistered tolerance.
\end{enumerate}
\subsection*{Baseline checklist}
For every key experiment, include: heuristic baseline, planner-off, and world-model-off outputs with identical seeds and scenario assignments.

\nocite{paper1,paper2,paper3,holm,efron}
\printbibliography

\end{document}
